\section{Where you are in every family}
\label{sec:torus}

Stand at a point $x$ on the pair of rulers. You are, say, three tenths of the
way from tick 40 to tick 41 of family A, and seven tenths of the way from
tick 36 to tick 37 of family B. Those two fractions, $0.3$ and $0.7$, are
where you are \emph{within the current member} of each family, and they are
all a pattern of coincidences can depend on: whether ticks coincide at $x$ is
a question about the fractions, not about the numbers 40 and 36.

\begin{definition}[The torus of counts]
\label{def:torus}
For $K$ families with counts $\cnt_1,\dots,\cnt_K$ the \emph{state} of the
superposition at $x$ is the point
\[
  \Phi(x)=\big(\cnt_1(x),\dots,\cnt_K(x)\big)\bmod 1,
\]
the fractional part of every count at once. For two families $\Phi(x)$ is a
point $(u_1,u_2)$ of the unit square. Walking off the right edge of the
square (family A completes a member) brings you back in at the left edge at
the same height, and walking off the top brings you back in at the bottom:
the square with its opposite edges glued is a \emph{torus}, written
$\tor^2$, and as $x$ moves, $\Phi(x)$ traces a path on it.
\end{definition}

\begin{lookup}{the torus as a square with its edges glued}
What to look up: \emph{why a square with opposite edges identified is called
a torus}. The one consequence this paper uses: a function on the torus is
the same thing as a function $I(u_1,u_2)$ of two ordinary variables that
repeats with period one in each variable separately. Nothing about the
torus's shape in three dimensions is ever used; only the gluing.
\end{lookup}

The state says where things are. What they look like is a separate thing,
and separating the two is the first idea worth having.

\begin{definition}[The picture on the torus]
\label{def:picture}
What the superposition \emph{looks like} at $x$ --- how much ink, how loud,
whether the detector fires --- depends on $x$ only through the state
$\Phi(x)$, by a fixed rule $I$ that knows nothing about $x$: the
\emph{picture on the torus}. The superposition is the composite
\[
  S \;=\; I\circ\Phi,\qquad S(x)=I\big(\Phi(x)\big).
\]
\end{definition}

The formula says: look up where you are in every family, then look up what
that state looks like. For two printed gratings $I(u_1,u_2)$ is ``ink if
$u_1$ is within the stroke of a line of A, or $u_2$ within a line of B'' ---
a picture on the torus made of two bands, one across and one up, and their
crossing. For two click trains into a coincidence detector $I(u_1,u_2)$ is
``fire if $u_1$ and $u_2$ are both within a click's width of zero'', a small
square at the corner of the torus. Everything about \emph{where} ---
pitches, spacings, the bend in a ruler, a drift in a clock --- lives in
$\Phi$. Everything about \emph{what} --- ink, click shape, brightness,
whether inks multiply or sounds add --- lives in $I$. Two superpositions
with the same $\Phi$ differ only in their picture, and two with the same
picture differ only in their state map. The rest of the paper is a list of
things that are properties of $\Phi$, things that are properties of $I$, and
things that are properties of the way an observer averages one over the
other.

\begin{figure}[t]
  \centering
  \begin{tikzpicture}[x=1cm, y=1cm, font=\small, >={Stealth[length=4pt]}]
    \begin{scope}
      \clip (0,0) rectangle (3.4,3.4);
      \foreach \i in {0,...,17} { \draw[ink, line width=0.45pt] ({\i*0.2},0) -- ({\i*0.2},3.4); }
      \foreach \i in {0,...,15} { \draw[ink, line width=0.45pt] ({\i*0.23+0.06},0) -- ({\i*0.23+0.06},3.4); }
    \end{scope}
    \draw[soft] (0,0) rectangle (3.4,3.4);
    \node[below=2pt] at (1.7,0) {the superposition $S=I\circ\Phi$};
    \draw[->, soft, thick] (3.9,1.7) -- (5.1,1.7) node[midway, above, ink] {$\Phi$};
    \begin{scope}[shift={(5.6,0)}]
      \clip (0,0) rectangle (3.4,3.4);
      \foreach \c in {0.4,0.8,...,6.8} { \draw[soft!45, thin] (\c,0) -- (0,\c); }
      \foreach \c in {-2.8,-2.1,...,2.8} { \draw[accent!70, thin] (\c,0) -- (\c+3.4,3.4); }
      \fill[ink!12] (0,0) rectangle (0.5,3.4);
      \fill[ink!12] (0,0) rectangle (3.4,0.5);
      % A neighbourhood's path: along the fast direction, whose slope is the
      % ratio of the two families' rates, 0.2/0.23 = 0.87.
      \draw[ink, line width=1.6pt] (1.1,0.35) -- (2.9,1.92);
      \fill[ink] (2.0,1.13) circle (1.6pt);
    \end{scope}
    \draw[ink] (5.6,0) rectangle (9,3.4);
    \node[below=2pt] at (7.3,0) {the torus $\tor^2$: $\cnt_1$ across, $\cnt_2$ up};
    \draw[accent, thin] (5.6,-0.75) -- (6.1,-0.75) node[right, ink, font=\scriptsize] {level lines of $\cnt_1-\cnt_2$ (slow)};
    \draw[soft!60, thin] (5.6,-1.05) -- (6.1,-1.05) node[right, ink, font=\scriptsize] {level lines of $\cnt_1+\cnt_2$ (fast)};
    \draw[ink, line width=1.6pt] (5.6,-1.35) -- (6.1,-1.35) node[right, ink, font=\scriptsize] {the path of a small neighbourhood of a point};
  \end{tikzpicture}
  \caption{\figlead{A superposition beside its torus.} Left: two families of
  lines a little apart in pitch. Right: the torus of their counts. The
  picture $I$ is the two shaded bands, a line of A across and a line of B
  up; ink at a point of the plane is $I$ looked up at that point's state.
  A small neighbourhood of a point maps to a short segment on the torus
  (black), pointing in the direction in which both counts advance --- the
  \emph{fast direction}. The level lines (contour lines) of $\cnt_1-\cnt_2$
  run nearly along the segment, so that combination barely changes across a
  neighbourhood; the level lines of $\cnt_1+\cnt_2$ are crossed many
  times. The first is slow and the second fast, and the whole theory is about
  what is slow.}
  \label{fig:torus}
\end{figure}

\subsection{Combined counts, and which of them are slow}
\label{sec:slow}

The difference $D=\cnt_A-\cnt_B$ of Figure~\ref{fig:rulers} is itself a
count: it rises by one per period of the bands, and the bands are its
members. So is the sum $\cnt_A+\cnt_B$, and so is any combination
\[
  k\cdot\cnt \;=\; k_1\cnt_1+k_2\cnt_2+\dots+k_K\cnt_K,
  \qquad k=(k_1,\dots,k_K)\ \text{whole numbers},
\]
a whole number of one count plus a whole number of another. The formula says:
combined counts are made of the counts by integer recipes, and the recipe is
the vector $k$. We write $k=(1,-1)$ for the difference, $(1,1)$ for the sum,
$(2,-1)$ for twice the first count minus the second. A combined count is a
count in its own right --- it runs upward like $D$ in Figure~\ref{fig:rulers}
--- and its fractional part depends only on the state $\Phi(x)$, because a
whole-number combination of whole numbers is a whole number. Every recipe is
a candidate pattern, all of them at once. Which ones show is a matter of how
fast each changes.

Here is why recipes with entries bigger than one matter, in one sentence
that \S\ref{sec:whosees} unpacks. A family with sharp edges, seen as a
function of its own count, is a sum of waves that repeat once, twice, three
times per member --- its \emph{harmonics} --- and the second harmonic of
family A against the first of family B makes a pattern of coincidences with
recipe $(2,-1)$, as surely as the first against the first makes $(1,-1)$.
A pair of rulers with ten and nineteen divisions to the inch shows it:
twice the first count minus the second rises by one per inch, and that is
the pattern you see, faint but there. Nineteen to ten is nearly two to one,
the ratio of an octave in music, so we call $(2,-1)$ the \emph{octave}
recipe.

\begin{definition}[Slow and fast; the heterodyne ratio]
\label{def:eta}
At a point $x$ each count has a local rate, its derivative $\nabla\cnt_i(x)$
--- a number on a line, a vector in the plane, giving members per unit of
distance. The rate of a combined count is $\nabla(k\cdot\cnt)=\sum_ik_i\nabla\cnt_i$.
For a recipe $k=(a,b)$ on two families, its \emph{heterodyne ratio} is
\[
  \het_{a,b}(x)\;=\;\frac{\lvert a\,\nabla\cnt_1+b\,\nabla\cnt_2\rvert}
                        {\tfrac12\big(\lvert a\,\nabla\cnt_1\rvert+\lvert b\,\nabla\cnt_2\rvert\big)},
\]
the rate of the combination divided by the average rate of the two
harmonics that form it. Its reciprocal is the number of members per band.
A recipe is \emph{slow} when $\het$ is small; we call $\het<\tfrac14$, more
than four members per band, the \emph{beat regime}, and it is where every
pattern in this paper lives.
\end{definition}

The word is the radio engineer's: to \emph{heterodyne} is to mix two rates
and keep their difference, and $\het$ says how clean the mix is. On the
rulers of Figure~\ref{fig:rulers}, with rates $10$ and $9$ per inch, the
difference has $\het=1/9.5$, nine and a half members per band, and the sum
has $\het=19/9.5=2$: half a member per band, which is no band at all. The
click trains of \S\ref{sec:coincidence} have rates $100$ and $101$ per
second and $\het=1/100.5$ for the difference; a hundred clicks per pulse.
Turn one ruler around so that its count runs the other way, and the sum
becomes the slow one; nothing in the theory prefers a difference. The
threshold $\tfrac14$ is where a pattern of coincidences stops being one:
with more than four members per band the bands are wider than the members
that make them, and with fewer there is no ``band'', only members. The
ratio is a number at every point of the plane, so it can be drawn as a map,
and the tool of \S\ref{sec:instrument} draws it: bright where no pattern can
form, dark where one must. From here on the members themselves --- the
ticks, the clicks, the lines, the fine texture a pooling observer loses ---
are the \emph{carrier}, radio's word again, and a slow recipe is a
\emph{beat}.

\begin{table}[t]
  \centering\small
  \begin{tabular}{@{}lll@{}}
    \toprule
    Symbol & Read it as & Introduced \\
    \midrule
    $\cnt_i(x)$ & the count of family $i$ up to $x$ & Definition~\ref{def:count} \\
    $\Phi(x)$ & the state: where you are in every family, on the torus $\tor^K$ & Definition~\ref{def:torus} \\
    $u=(u_1,u_2)$ & a point of the torus; $x$ is always a point of the line or plane & Definition~\ref{def:torus} \\
    $I$ & the picture on the torus: what a state looks like & Definition~\ref{def:picture} \\
    $S=I\circ\Phi$ & the superposition itself & Definition~\ref{def:picture} \\
    $k=(k_1,\dots,k_K)$ & a recipe: whole numbers & \S\ref{sec:slow} \\
    $k\cdot\cnt$ & the combined count with recipe $k$ & \S\ref{sec:slow} \\
    $D$ & the difference of two counts, recipe $(1,-1)$ & Figure~\ref{fig:rulers} \\
    $\het$ & the heterodyne ratio: one over the members per band & Definition~\ref{def:eta} \\
    $\hat I(k)$ & the strength of recipe $k$'s wave in the picture & \S\ref{sec:sweep} \\
    $w$, $\mathcal E_w$ & a slide (a recipe of whole numbers) and the average along it & Theorem~\ref{thm:sweep} \\
    $\E[\,I\mid\text{slow}\,]$ & the picture averaged over everything the slow recipes do not fix & \S\ref{sec:sweep} \\
    $T_I(D)$ & the tent: the average as a function of $D$ & Corollary~\ref{cor:cases} \\
    $W$, $\rho$, $\widehat W$ & an observer's window, its width, and its response to a wave & \S\ref{sec:window} \\
    $N$ & an observer's front end: what it does to each value before pooling & \S\ref{sec:observer} \\
    $O$ & an observer: a front end, then a window, then anything & Theorem~\ref{thm:universality} \\
    \bottomrule
  \end{tabular}
  \caption{\figlead{Every symbol in the paper}, and where it is introduced.
  The lower half is used from \S\ref{sec:pooling} on.}
  \label{tab:symbols}
\end{table}

\begin{tryit}{the map of where a beat can form}
In \Moire\ (\S\ref{sec:instrument}), make two layers of parallel lines a few
percent apart in pitch and turn on the \emph{ratio} view. The frame darkens
where $\het<\tfrac14$: the bands you see in the ordinary view sit exactly
inside the dark region, and as you change the pitches the dark region
predicts where the bands go before they get there. Rotate one layer: the
dark region moves to where the \emph{sum} is slow, and the bands follow.
\end{tryit}
